\documentclass[journal]{IEEEtran}
\usepackage{cite}
\usepackage{amsmath,amssymb}
\usepackage{graphicx}
\usepackage{booktabs}
\usepackage{longtable}
\usepackage{url}
\usepackage[hidelinks]{hyperref}

\graphicspath{{figures/}}

\begin{document}

\title{Causal Consistency as a Defense: Detecting and Auditing Provenance Graph Poisoning}

\author{\IEEEauthorblockN{N.~Manoj}
\IEEEauthorblockA{Independent Researcher\\
Email: nehamanoj1105@users.noreply.github.com}}

\maketitle

\begin{abstract}
Provenance-based intrusion detection consumes a stored graph and assumes it is
authentic. We examine what happens when that assumption fails, and whether the
causal invariants that a provenance graph must satisfy by construction are
sufficient to reveal post-collection tampering without any training. We assemble
fifteen deterministic rules over structural, temporal and semantic invariants
into a rule engine and evaluate it against a GraphSAGE baseline on controlled
synthetic poisoning and on real DARPA Transparent Computing Engagement~3 Theia
provenance. Over ten seeds, the rule engine reaches F1 $0.805\pm0.063$ on
synthetic poisoning, with a genuine continuous score giving ROC-AUC
$0.953\pm0.031$, while a leakage-free GraphSAGE reaches F1 $0.322\pm0.194$; the
paired difference is $+0.482$ (Wilcoxon $p=0.002$). The rule engine's accuracy is
unchanged as invariant-preserving mimicry grows from none to heavy. On real
Theia graphs with synthetic post-collection poisoning it reaches F1
$0.835\pm0.046$ and $0.850\pm0.060$, where GraphSAGE is near chance. We give
particularly close attention to evaluation design: the conventional protocol of
training and thresholding a baseline on the same edges inflates GraphSAGE F1 by
$0.342$ absolute, and a reported ROC-AUC of $1.000$ for the rule engine is not a
computed value. We separate real from synthetic evaluation throughout and release
the complete corrected protocol, per-seed results and verification scripts.
\end{abstract}

\begin{IEEEkeywords}
provenance, intrusion detection, graph poisoning, causal consistency,
GraphSAGE, evaluation methodology, reproducibility, DARPA Transparent Computing
\end{IEEEkeywords}

\section{Introduction}
The provenance of a computer system is its causal history: which process read
which file, which process spawned which child, which process opened which
network connection. Captured faithfully, it is a directed graph rich enough to
reconstruct an intrusion after the fact and, increasingly, to detect it as it
happens. Systems such as Unicorn~\cite{han2020unicorn}, NoDoze~\cite{hasan2021nodoze},
HOLMES~\cite{milajerdi2019holmes}, Poirot~\cite{peirce2021poirot},
FLASH~\cite{chen2021flash} and MAGIC~\cite{li2021magnet} build on the assumption
that the provenance graph is a trustworthy record.

That assumption holds only up to the point where the graph stops being produced
and starts being stored. Once events are serialized, shipped and indexed, the
graph becomes data, and data can be edited by an adversary with access to the
storage. Deleting the events that record a privilege escalation, inserting
events that suggest a process behaved benignly, reordering events to break a
causal chain, or re-attributing an action to an innocent actor are all plausible
operations against a stored graph. A downstream detector then reasons over a
fiction with no indication that anything is wrong.

Our question is whether causal consistency is enough to catch this. The events
in a provenance graph are not arbitrary; they must satisfy invariants that come
from what the events mean. An execute edge must point at a file. A read edge
must point at a file. A child process cannot be spawned before its parent
exists. Event sequence numbers within a process must increase. A process that
appears to act must at some point have been spawned. These invariants are cheap
to check and require no training data, which raises the question of how far they
take us against an adversary who edits the graph.

We answer with an empirical study and, deliberately, a methodological one. We
build a fifteen-rule engine from these invariants and compare it to GraphSAGE, a
representative learned detector, on the same poisoned instances. The comparison
favors the rule engine, but the more useful result is why: a learned baseline
evaluated in the conventional way---fit on all edges, threshold chosen to
maximize F1 on the scored labels---has most of its reported advantage supplied
by that choice. When the split is made honest, the advantage largely disappears.
Because this failure mode is easy to reproduce and easy to miss, we treat it as
a first-class result rather than a footnote.

\subsection{Contributions}
\begin{itemize}
\item A deterministic, training-free detector for post-collection provenance
tampering, built from fifteen structural, temporal and semantic invariants
(Section~\ref{sec:framework}).
\item A rigorous evaluation protocol with disjoint train, validation and test
edges, validation-only threshold selection, and test-set isolation from weights,
threshold, hyperparameters and attack intensity (Section~\ref{sec:protocol}).
\item A ten-seed comparison on synthetic poisoning and an evaluation on two real
Theia scenarios with synthetic post-collection poisoning, with paired
significance tests (Sections~\ref{sec:main}--\ref{sec:robust}).
\item A genuine continuous anomaly score for the rule engine, replacing a
reported ROC-AUC of $1.000$ that was a fallback constant rather than a
computation; the computed value is $0.953\pm0.031$ (Section~\ref{sec:rocauaudit}).
\item A leakage analysis quantifying a $0.342$ absolute F1 inflation for the
conventional GraphSAGE protocol (Section~\ref{sec:integrity}).
\item A complete release of protocols, raw per-seed results and verification
scripts so that every number can be re-derived (Section~\ref{sec:repro}).
\end{itemize}

\section{Background and Related Work}
\subsection{Provenance-based detection}
Specification-based systems (SLEUTH~\cite{hossain2017sleuth},
HOLMES~\cite{milajerdi2019holmes}, Poirot~\cite{peirce2021poirot}) encode expert
knowledge of attack behavior and are precise but require continued maintenance
and do not adapt to novel tactics. Statistical systems score the rarity of
events or paths (NoDoze~\cite{hasan2021nodoze}, ProvDetector), and are prone to
flagging unusual-but-benign activity. Learning-based systems learn
representations of benign behavior: Unicorn~\cite{han2020unicorn} at coarse
granularity, FLASH~\cite{chen2021flash} and MAGIC~\cite{li2021magnet} at finer
granularity, with GNN backbones such as GraphSAGE~\cite{hamilton2017graphsage}
and GAT.

\subsection{Adversarial robustness of graph detectors}
Graph-based detectors have been shown vulnerable to mimicry, in which an
adversary interleaves benign activity with malicious activity so that local
aggregation dilutes the anomaly signal~\cite{jiang2023mimicry}. Defenses have
moved toward dynamic neighbor selection and reinforcement-learned aggregation.
Our setting differs: the adversary is not trying to look benign \emph{inside} the
graph, but to edit the graph itself. A detector that verifies invariants is
therefore addressing a different question than a detector that classifies
behavior.

\subsection{Provenance integrity}
The integrity problem for provenance has mostly been approached through
capture-side mechanisms---trusted logging, tamper-evident storage, cryptographic
chaining. Our work is complementary: given a graph whose custody is uncertain, we
ask how much can be recovered from internal consistency alone, as a lightweight
layer before any cryptographic mechanism is available.

\section{Problem Formulation}
Let $G=(V,E)$ be a provenance graph produced by a capture pipeline and then
stored. We consider an adversary that modifies $E$ between capture and analysis.
The analysis receives a graph $G'$ and must decide, for each edge, whether it
witnesses tampering. Formally, we seek a detector $D:G'\rightarrow\{0,1\}^{|E'|}$
and, ideally, a score $s:G'\rightarrow\mathbb{R}^{|E'|}$.

The ground truth is the set of edges the adversary changed. This is well defined
for insertion, reordering and forgery, because those operations leave a
(modified) edge in $G'$. It is not well defined for deletion, because the
deleted edge is, by construction, absent from $G'$; an edge-level detector
cannot mark an edge that is not there. We therefore do not treat deleted edges
as scored positives. Their detectability comes only from breaking a higher-level
invariant elsewhere---a parent process left without a spawn record, a gap in a
sequence---and we report deletion separately.

\section{Threat Model}
We assume the attacker can read and write the stored graph, but cannot modify
the detector, disable the parser, or compromise the kernel that produced the
records (Table~\ref{tab:threat}). The attacker may delete, insert, reorder and
forge edges, and may edit node and edge attributes. The attacker is gray-box:
the schema and the general class of detector are known, but not detector
parameters or learned weights. The goal is to hide malicious activity while
preserving as much observable structure as possible. Timestamp manipulation,
attribute forgery and targeted stealth poisoning are in scope as future attack
classes but are not evaluated here.

\begin{table}[t]
\centering
\small
\caption{Threat model summary.}
\label{tab:threat}
\begin{tabular}{p{3.0cm}p{5.0cm}}
\toprule
Aspect & Assumption \\
\midrule
Trusted & kernel, audit logging, parser, detector \\
Untrusted & stored graph, intermediate representations, storage \\
Attacker can & delete, insert, reorder, forge edges; edit attributes \\
Attacker cannot & modify detector/parser, compromise kernel \\
Knowledge & schema + detector class (gray-box) \\
Goal & hide malicious activity, preserve structure \\
\bottomrule
\end{tabular}
\end{table}

\section{Provenance Representation}
A provenance graph $G=(V,E)$ has typed nodes and edges. Node types are process,
file, user and network endpoint. Edge types are read, write, execute, connect,
spawn and delete. Each edge carries source and target identifiers, a type and a
timestamp; each event carries a sequence number within its owning process. This
is the common data model of the DARPA Transparent Computing program, whose
records are serialized in Avro; our parser decodes that schema into the
representation above. The rules below are stated entirely in terms of these
types, the identifiers and the timestamps, so they apply to any capture that
produces this model.

\section{Proposed Detection Framework}
\label{sec:framework}
The detector is a rule engine that evaluates fifteen invariants over the graph
and returns, for each, the set of violations, each attributed to an edge or a
node. We group the rules by the kind of invariant they check.

\subsection{Structural rules}
Well-formedness of types and topology: no duplicate edge tuples; no duplicate
event identifiers; spawn edges must target process nodes; execute, read, write
and delete edges must target file nodes; connect edges must target network
endpoints; no self-referential edges; no edges referencing nodes that are not
present; and no active process lacking a parent spawn edge.

\subsection{Temporal rules}
Ordering invariants: non-negative timestamps; a child's spawn must not precede
its parent's spawn; process activity must not precede the process spawn; event
sequence numbers must be contiguous enough to reveal gaps; and timestamps must
be monotone within a process stream.

\subsection{Semantic rules}
Lineage invariants: a process that acts without any observed parent spawn is
flagged, capturing attempted laundering of provenance through re-attribution.

\subsection{Continuous anomaly score}
For ranking and for AUC, every violation contributes a severity weight (high
$3$, medium $2$, low $1$) to the edge it is attributed to, either directly by
edge identifier or, for node-attributed violations, to the incident edges. The
binary decision is score greater than zero. This score is a genuine, monotone
ranking signal derived from the violation structure; we use it, and only it,
when reporting ROC-AUC and PR-AUC. No AUC in this paper is produced by a
fallback constant.

\section{GraphSAGE Baseline}
\label{sec:graphsage}
The learned baseline is a two-layer GraphSAGE encoder~\cite{hamilton2017graphsage}
with an edge classifier. Nodes are featurized by a one-hot node type and degree;
edges by a one-hot edge type and a normalized timestamp. The baseline exists to
test whether learning can exceed fixed invariants under a fair protocol. It is
not tuned to win, its feature set is deliberately minimal, and Section
\ref{sec:integrity} shows that its apparent strength in the original repository
came largely from the evaluation protocol rather than from the model.

\section{Experimental Methodology}
\label{sec:protocol}
\subsection{Datasets}
\label{sec:datasets}
Two data sources are used. The first is a synthetic provenance generator
producing a graph of thirty processes, forty files and ten network entities,
which provides exact ground truth and allows the attack intensity to be varied
at will. The second is real provenance from DARPA Transparent Computing
Engagement~3, the Theia streams~\cite{darpa2020tc}. Four benchmark scenarios are
named in the original material (1r, 3, 5m, 6r). Two were obtained and parsed:
scenario~3 with 20{,}157 nodes and 297{,}777 edges, and scenario~5m with 34{,}835
nodes and 464{,}858 edges. Both parse counts match the documented statistics
exactly. Scenarios~1r and~6r could not be downloaded, because the public mirror
enforced a download quota (HTTP~403); we report the two obtained scenarios and
state the missing coverage as a limitation rather than estimating it
(Table~\ref{tab:datasets}).

Real-data experiments are run on the first 50{,}000 edges of each real graph,
matching the loader's behavior; the truncation is explicit. We emphasize the
category separation used throughout: the real-data experiments are \emph{real
provenance with synthetic post-collection poisoning}. They are not real attack
labels, and we never describe them as detection of naturally occurring attacks.

\subsection{Poisoning operators}
\label{sec:poisoning}
For each seed we inject five deletions, five insertions, five reorderings and
five dependency forgeries, for twenty events. The injector operates on a deep
copy of the clean graph, records for each event its original state, its modified
state, an expected detection target and a flag indicating whether the operation
actually changed the graph. It then asserts that the clean graph is unmodified.
This matters because a naive implementation can mutate its input in place, and
because some reorderings can be no-ops if they do not change the semantic order;
we mark such events ineffective and exclude them from detectable positives.

Twenty events yield fifteen scored positives, since the five deletions leave no
final-graph edge. Detection metrics are computed over the final graph
(\emph{final-graph universe}). An alternative convention, which adds the
identifiers of deleted edges to the scoring universe, raises F1 from $0.805$ to
$0.845$; this is quantified in Section~\ref{sec:poisongt} and we use the
final-graph universe in all reported results.

\subsection{Leakage-free protocol}
\label{sec:leakprotocol}
For the learned baseline, edges are partitioned into disjoint training,
validation and test sets. Weights are fit on the training edges only; the
decision threshold is selected on the validation edges only; and every reported
metric is computed on the untouched test edges. Message passing for test-node
features uses training and validation edges only, so test edges influence
neither model weights nor threshold. Neither hyperparameters nor attack
intensity is tuned on the test set. For the training-free rule engine there is
no split to leak, which is precisely why the leakage below affects only the
baseline.

\subsection{Metrics}
We report precision, recall, F1, accuracy, balanced accuracy, MCC, FPR, FNR,
specificity, ROC-AUC and PR-AUC. ROC-AUC and PR-AUC are reported only where a
genuine continuous score exists and both classes are present. For the rule
engine this is the severity-weighted score of Section~\ref{sec:framework}; for
GraphSAGE it is the classifier probability.

\subsection{Seeds and statistics}
Synthetic, mimicry, ablation and real-Theia experiments use ten seeds
($1,7,13,21,42,99,123,256,512,1024$); generalization uses five
($1,7,13,21,42$). Each seed produces a freshly generated graph and a fresh set
of attack instances. We report the mean, sample standard deviation, median,
minimum, maximum and a 95\% confidence interval (Student-$t$, $n-1$ degrees of
freedom). Paired comparisons between detectors use the Wilcoxon signed-rank test
on the same test instances, with Cohen's $d_z$ as an effect size.

\section{Data Parsing and Preprocessing}
Raw Theia records are Avro-serialized common data model entries nested under a
datum envelope. The parser walks each record, extracts the typed subject and
predicate, and emits an edge with source, target, type and timestamp. Nodes are
materialized on demand from identifiers. Parsing the two obtained scenarios
yields zero skipped records, and the counts agree with the published statistics,
which is our validation that the decode is correct. No attack labels are read
from the real data; the DARPA ground-truth report was obtained for context only
and is not used as a detection target.

\section{Evaluation Protocol}
\label{sec:evaluation}
Detection is evaluated at the edge level against the injected ground truth,
under the final-graph universe. For the rule engine, the binary decision is
score greater than zero and the continuous score is used for AUC. For GraphSAGE
the protocol is the leakage-free one of Section~\ref{sec:leakprotocol}. Real and
synthetic evaluations are never pooled: Table~\ref{tab:main} is synthetic,
Table~\ref{tab:realpoison} is real-provenance-with-synthetic-poisoning, and the
two are reported and discussed separately.

\section{Main Detection Results}
\label{sec:main}
Table~\ref{tab:main} reports the ten-seed synthetic comparison. The rule engine
reaches F1 $0.805\pm0.063$, with precision $0.715\pm0.095$, recall
$0.933\pm0.063$ and specificity $0.964\pm0.017$; MCC is $0.795\pm0.064$. The
leakage-free GraphSAGE reaches F1 $0.322\pm0.194$, precision $0.342\pm0.309$ and
recall $0.502\pm0.364$. The learned baseline is not competitive once evaluated
honestly, and its seed-to-seed variance is roughly three times larger. Under a
paired test over the same instances, the F1 difference is $+0.482$ with a 95\%
confidence interval of $[0.328, 0.636]$, Cohen's $d_z=2.24$ and a Wilcoxon
$p=0.002$; the ROC-AUC difference ($+0.124$, $p=0.105$) is not statistically
distinguished at this sample size (Table~\ref{tab:significance}). We report that
non-result rather than claiming a uniform advantage.

Figure~\ref{fig:main} shows the per-metric comparison. The rule engine's high
recall and specificity are the expected signature of a detector that checks
invariants: it catches almost all injected structure while rarely flagging
untouched edges. Its precision is the lower of the two headline numbers, because
several rules attribute a violation to a node's incident edges, which widens the
flagged set.

\section{Cross-Dataset Generalization}
\label{sec:generalization}
Because the rule engine has no fitted parameters, it can be applied to a new
graph unchanged. Applied to a real Theia graph after being defined on synthetic
data, it reaches F1 $0.857\pm0.029$; moving between the two real scenarios it
reaches $0.837\pm0.021$ (Table~\ref{tab:generalization}). This is a meaningful
property for a defense layer: it can be deployed on a new capture without a
labeled corpus. We deliberately make no corresponding claim for GraphSAGE. A
genuine transfer across a shared feature space is not implemented in this study,
so the GraphSAGE column in Table~\ref{tab:generalization} is a target-trained
model, and we label it as such rather than presenting it as generalization.

\section{Robustness to Mimicry}
\label{sec:robust}
Mimicry is an adversary's attempt to hide tampering inside plausible noise. We
generate camouflage edges that preserve the degree, type and timestamp
distributions of the surrounding graph and that satisfy every implemented
invariant, and increase their number from none to heavy (Table~\ref{tab:mimicry},
Figure~\ref{fig:mimicry}). The rule engine's F1 is $0.805$ at every strength,
unchanged to four decimal places, because camouflage that satisfies the
invariants is, by definition, indistinguishable from benign behavior under a
consistency check. GraphSAGE ranges from F1 $0.304$ to $0.542$ across strengths
and remains well below the rule engine.

Two specification details are essential to interpret this correctly. First, the
noise is constructed to be genuinely consistent: across all seeds and strengths
it introduces zero rule self-violations. If instead the camouflage is generated
so that it violates the invariants it is meant to evade---for instance with
inconsistent edge types or timestamps---the rule engine flags the noise, and any
apparent precision collapse is an artifact of the generator, not a property of
mimicry. Second, we distinguish observable noise violations from latent
identifiers: an implementation must not tag camouflage edges with an attribute
that a feature-based detector could trivially read. We verify that no detector
in this study consumes such an attribute.

\section{Poisoning Experiments and Ground-Truth Validity}
\label{sec:poisongt}
Table~\ref{tab:realpoison} reports the real-provenance experiments. On scenario~3
the rule engine reaches F1 $0.835\pm0.046$ and on scenario~5m F1
$0.850\pm0.060$; GraphSAGE reaches F1 $0.034\pm0.105$ and $0.010\pm0.009$. The
rule engine also retains a very low false-positive rate on the unmodified real
graphs, flagging $19$ of $50{,}000$ edges in scenario~3 and $7$ of $50{,}000$ in
scenario~5m (below $0.04\%$).

The ground-truth audit is worth stating explicitly. Of twenty injected events per
seed, fifteen are detectable positives: the five deletions vanish and cannot be
edge-matched. The remaining operations are real modifications of final-graph
edges. Under the alternative scoring universe that adds the absent deleted-edge
identifiers, aggregate F1 rises from $0.805$ to $0.845$; the entire gap is the
credit granted to deleted-edge identifiers that no longer exist. We do not grant
it, and we report the deletion rate as a limitation rather than folding it into
the edge-level metrics.

\section{Rule Ablation}
\label{sec:ablation}
Table~\ref{tab:ablation} reports rule-group removal and leave-one-rule-out over
ten seeds. The structural group carries the detector: removing it drops F1 from
$0.805$ to $0.328$. Removing the temporal group costs little ($0.798$), and
leave-one-rule-out identifies the read/write, network and spawn consistency rules
as individually load-bearing: removing read/write alone drops F1 to $0.517$,
removing spawn to $0.742$, removing network to $0.746$. Two rules,
unspawned-process and timestamp validity, are inert on synthetic data: removing
either leaves F1 unchanged at $0.805$. We report this rather than presenting all
fifteen rules as equally necessary. Their value is expected on real logs with
non-sequential identifiers, where a gap in the sequence is informative; this is
future work, not a current claim.

\section{GraphSAGE Leakage Analysis and Evaluation Integrity}
\label{sec:integrity}
The most consequential finding of this study is about evaluation, not modeling.
The conventional protocol for the learned baseline trains on all edges and then
selects the decision threshold that maximizes F1 on the same labels it reports.
On the identical seed, this protocol yields GraphSAGE F1 $0.588$, while the
leakage-free protocol yields $0.246$ (Figure~\ref{fig:leakage}). The difference,
$0.342$ absolute, is not model quality; it is the threshold being chosen with
knowledge of the test answers.

Three specific leakage channels are present in such a protocol and absent from
ours. First, training edges are also test edges, so the model sees the instances
it is scored on. Second, the threshold is optimized on the scored set, which
directly maximizes the reported metric. Third, message passing over the full
graph lets test edges influence the features used for test nodes. Removing all
three is what separates the two numbers above. Because the rule engine has no
fitted parameters and no threshold search, it is invariant to this class of
error, which is a practical advantage of a training-free detector beyond
accuracy.

We also audit a reported ROC-AUC of $1.000$ for the rule engine. It is not a
computed quantity: the evaluation record for the rule engine contains no ROC-AUC,
and the reported value is a fallback constant applied when the field is missing.
Such a value is meaningless and we do not use it. The genuine AUC, computed from
the severity-weighted score of Section~\ref{sec:framework}, is
$0.953\pm0.031$ ($\text{PR-AUC}=0.740\pm0.080$), reported in
Table~\ref{tab:main} and shown for a representative seed in
Figure~\ref{fig:rocpr}.

\section{ROC-AUC Verification and Continuous Scoring}
\label{sec:rocauaudit}
We verify the AUC claim independently. The rule engine emits binary violations;
to obtain a ranking score we aggregate violations per edge with severity weights
and treat the resulting non-negative score as a ranking signal. This produces a
genuine ROC curve and a genuine precision--recall curve
(Figure~\ref{fig:rocpr}), with edge-level AUC computed by the standard estimator.
The value is $0.953$ on the seed shown and $0.953\pm0.031$ across ten seeds, i.e.
strong but not perfect. We report PR-AUC (average precision) alongside, since the
positive class is a small minority and average precision is the more informative
summary under imbalance. We do not report an AUC for the rule engine anywhere
that a legitimate continuous score is not computed.

\section{Scalability and Computational Cost}
\label{sec:scalability}
Table~\ref{tab:scalability} and Figure~\ref{fig:scalability} report runtime,
memory and throughput from ten thousand to one million edges, five repetitions
per size. We time construction, poisoning, rule-engine inference and end-to-end
execution separately, and we report GraphSAGE training and inference separately
from rule-engine inference, because comparing a training cost to an inference
cost is not meaningful.

Construction and inference are both linear in the edge count. Inference takes
$0.37$~s at ten thousand edges and $34.5$~s at one million; construction takes
$0.18$~s and $21.2$~s. Peak Python allocation grows from about $37$~MB to
$971$~MB, close to one kilobyte per edge, and throughput is stable near
$2.9\times10^{4}$ edges per second across the full range. The flat throughput is
the expected result of a single linear pass over the events and indicates no
super-linear behavior hiding in the rule set. Poisoning is the dominant cost,
roughly sixty percent of end-to-end, because the corrected injector deep-copies
the graph and checks integrity on every event.

\section{Statistical Analysis}
\label{sec:statistics}
Table~\ref{tab:significance} reports the paired analysis. For F1, the mean paired
difference is $+0.482$ with a 95\% confidence interval of $[0.328, 0.636]$,
Cohen's $d_z=2.24$ and Wilcoxon $p=0.002$ ($n=10$), a large and statistically
supported effect. For ROC-AUC the mean difference is $+0.124$ with an interval of
$[0.005, 0.242]$ and $p=0.105$: at this sample size the AUC difference is not
distinguished from zero, and we say so rather than reporting only the significant
metric. All tests use paired observations on identical test instances, so the
comparison does not depend on separately drawn evaluation sets.

\section{Discussion}
The results support a specific, bounded claim: post-collection tampering that
breaks a provenance graph's internal causal invariants is detectable, in a
training-free way, on both synthetic and real graph structure, with linear cost
and strong robustness to invariant-preserving noise. The results also support a
methodological claim that is easy to overlook: a learned baseline can appear far
stronger than it is when its threshold is chosen on the evaluation labels, and
removing that leak can reverse a comparison. Both claims are useful to the
community---the first for defense, the second for how defense papers are
evaluated.

They also mark the boundary of the contribution. The rules check
well-formedness, not intent: an adversary who edits the graph while preserving
all invariants is not caught by this layer, and no statistical detector would
catch a modification that makes the graph \emph{more} consistent. The rule engine
is best understood as a cheap, deterministic, auditable first pass that a
heavier detector or a cryptographic integrity mechanism can build on.

\section{Limitations}
The evaluation covers two of four Theia scenarios; scenarios 1r and 6r were
unavailable and no results are estimated for them. All DARPA evaluations use
synthetic post-collection poisoning rather than real attacks, so they measure
detection of injected tampering on real graph structure, not detection of real
intrusions. Deletion is only partially detectable: a deleted edge cannot be
scored by edge matching, so we exclude it and do not claim a deletion detection
rate. Real graphs are truncated to their first 50{,}000 edges. The GraphSAGE
feature set is minimal and no shared-feature cross-graph transfer is
implemented, so the generalization study covers the rule engine only. Finally,
the class balance on real graphs is extreme (fifteen positives in fifty thousand
edges), so precision there is sensitive to a handful of false positives.

\section{Reproducibility}
\label{sec:repro}
We release the corrected protocol, the poisoning and mimicry generators, the
leakage-free harness, the raw per-seed result files, the aggregation scripts,
and the verification scripts. Every headline number is re-derived from the raw
per-seed files by an independent checker; the check recomputes each metric from
its confusion matrix (1{,}890 metric recomputations, zero mismatches) and
cross-checks 226 headline values against the tables (zero discrepancies). The
pipeline is deterministic: re-running seeds 1, 42 and 1024 in memory reproduces
the stored raw files exactly. The repository's original results are preserved
unmodified, and every deviation between the original protocol and the corrected
one is documented with its cause. Table~\ref{tab:repro} summarizes the release
contents.

\begin{table}[t]
\centering
\small
\caption{Reproducibility artifacts.}
\label{tab:repro}
\begin{tabular}{p{3.0cm}p{4.7cm}}
\toprule
Artifact & Contents \\
\midrule
Raw results & per-seed JSON for every experiment \\
Aggregation & scripts that build all tables from raw results \\
Validation & independent recomputation + cross-check scripts \\
Protocols & poisoning, mimicry, leakage-free harness \\
Manifests & dataset provenance, checksums, parser \\
\bottomrule
\end{tabular}
\end{table}

\section{Conclusion}
We asked whether the causal invariants of a provenance graph are enough to reveal
post-collection poisoning, and found that a small deterministic rule set exploits
them effectively. Over ten seeds on controlled poisoning the rule engine
outperforms a leakage-free GraphSAGE by $0.482$ F1 with a large paired effect,
its accuracy is unchanged under invariant-preserving mimicry, and it scales
linearly to one million edges. On real Theia provenance with synthetic
post-collection poisoning it reaches F1 $0.83$--$0.85$ where GraphSAGE is near
chance. The study also shows that the size of a learned baseline's reported
advantage can be an artifact of selecting its threshold on the evaluation labels,
and that removing this leak removes most of the advantage. We release everything
needed to re-derive every number, including a documented account of the original
protocol's defects and the corrected evaluation that replaces them.

\input{tables/table3_main_detection}
\input{tables/table9_real_data_poisoning}
\input{tables/table6_mimicry}
\input{tables/table7_rule_ablation}
\input{tables/table5_cross_dataset}
\input{tables/table8_scalability}
\input{tables/table4_statistical_comparison}
\input{tables/table1_dataset_statistics}

\begin{figure}[t]
\centering
\includegraphics[width=0.95\columnwidth]{fig_main_performance.pdf}
\caption{Synthetic controlled poisoning, ten seeds: Rule Engine versus
leakage-free GraphSAGE (mean and 95\% CI).}
\label{fig:main}
\end{figure}

\begin{figure}[t]
\centering
\includegraphics[width=0.95\columnwidth]{fig_roc_pr.pdf}
\caption{Rule Engine ROC and precision--recall curves from a genuine continuous
anomaly score (synthetic seed 42).}
\label{fig:rocpr}
\end{figure}

\begin{figure}[t]
\centering
\includegraphics[width=0.8\columnwidth]{fig_leakage.pdf}
\caption{Effect of evaluation leakage on GraphSAGE F1 for the same seed.}
\label{fig:leakage}
\end{figure}

\begin{figure}[t]
\centering
\includegraphics[width=0.95\columnwidth]{fig_mimicry.pdf}
\caption{Robustness to invariant-preserving mimicry (ten seeds).}
\label{fig:mimicry}
\end{figure}

\begin{figure}[t]
\centering
\includegraphics[width=0.95\columnwidth]{fig_ablation.pdf}
\caption{Rule ablation and leave-one-rule-out (ten seeds).}
\label{fig:ablation}
\end{figure}

\begin{figure}[t]
\centering
\includegraphics[width=0.95\columnwidth]{fig_scalability.pdf}
\caption{Scalability, five repetitions per size.}
\label{fig:scalability}
\end{figure}

\begin{figure}[t]
\centering
\includegraphics[width=0.95\columnwidth]{fig_generalization.pdf}
\caption{Cross-dataset generalization (five seeds). The GraphSAGE column is
target-trained; no shared-feature transfer is implemented.}
\label{fig:generalization}
\end{figure}

\begin{thebibliography}{99}
\bibitem{han2020unicorn} X.~Han, T.~Pasquier, A.~Bates, J.~Mickens, and
M.~Seltzer, ``Unicorn: Runtime provenance-based detector for advanced
persistent threats,'' in \emph{NDSS}, 2020.
\bibitem{hsieh2018kobold} K.~Hsieh \emph{et al.}, ``Kobold: Evaluating
decentralized access control,'' in \emph{IEEE S\&P Workshops}, 2018.
\bibitem{jiang2023mimicry} P.~Jiang \emph{et al.}, ``Mimicry attacks against
provenance graph based intrusion detection systems,'' in \emph{NDSS}, 2023.
\bibitem{wang2020you} Q.~Wang \emph{et al.}, ``You are what you do: Hunting
stealthy malware via data provenance analysis,'' in \emph{NDSS}, 2020.
\bibitem{li2021magnet} Z.~Jia, Y.~Xiong, Y.~Nan, Y.~Zhang, J.~Zhao, and
M.~Wen, ``MAGIC: Detecting advanced persistent threats via masked graph
representation learning,'' in \emph{USENIX Security}, 2024.
\bibitem{chen2021flash} M.~A.~Rehman and S.~Ahmadi, ``FLASH: A comprehensive
approach to intrusion detection via provenance graph representation learning,''
in \emph{IEEE S\&P}, 2022.
\bibitem{king2003backtracking} S.~T.~King and P.~M.~Chen, ``Backtracking
intrusions,'' in \emph{ACM SOSP}, 2003.
\bibitem{hossain2017sleuth} M.~N.~Hossain \emph{et al.}, ``SLEUTH: Real-time
attack scenario reconstruction from COTS audit data,'' in \emph{USENIX
Security}, 2017.
\bibitem{milajerdi2019holmes} S.~M.~Milajerdi, R.~Gjomemo, B.~Eshete,
R.~Sekar, and V.~N.~Venkatakrishnan, ``HOLMES: Real-time APT detection through
correlation of suspicious information flows,'' in \emph{IEEE S\&P}, 2019.
\bibitem{peirce2021poirot} S.~M.~Milajerdi, B.~Eshete, R.~Gjomemo, and
V.~N.~Venkatakrishnan, ``Poirot: Aligning attack behavior with kernel audit
records for cyber threat hunting,'' in \emph{ACM CCS}, 2019.
\bibitem{hasan2021nodoze} W.~U.~Hassan, M.~A.~Noureddine, P.~Datta, and
A.~Bates, ``NoDoze: Combatting threat alert fatigue with automated provenance
triage,'' in \emph{NDSS}, 2019.
\bibitem{hamilton2017graphsage} W.~L.~Hamilton, R.~Ying, and J.~Leskovec,
``Inductive representation learning on large graphs,'' in \emph{NeurIPS}, 2017.
\bibitem{darpa2020tc} DARPA I2O, ``Transparent Computing Engagement 3 data
release,'' 2020. [Online]. Available:
\url{https://github.com/darpa-i2o/Transparent-Computing}
\bibitem{pasquier2018practical} T.~Pasquier \emph{et al.}, ``Practical
whole-system provenance capture,'' in \emph{ACM SoCC}, 2017.
\end{thebibliography}

\end{document}
